\section{Transformer DC Extension Results}
\label{sec:vit-extension}
\appneedspace{8\baselineskip}
\begin{center}
    \captionsetup{hypcap=false}
    \captionof{table}[ViT-B/16: extension results with confidence intervals]{ViT-B/16 extension results on the same ImageNet-S test split (565 images; MaxS/AvgS: the fixed 200-image subset), with 95\% image-bootstrap confidence intervals. Gradient is copied from Table~\ref{ptnb:tab:test-ci-vit}; the other configurations were selected on the same 50 validation images. Labels loc/faith denote the unchanged selection criteria. Bold/dark shading marks the best row; $^{\approx}$/light shading marks no significant difference under Holm adjustment across all row pairs within each metric of this table (Section~\ref{ptnb:app:stats}).}
    \label{tab:vit-extension}
\end{center}
\captionsetup{hypcap=true}
\appinputgenci{tables/vit_extension/generated/extension_ci_longtable}
\noindent $B$ is the one-based block index (LN1 readout). SplitCAM uses the signed paper formulas without ReLU; $F$, $L$ and $U$ denote fixed attention and LayerNorm, fixed LayerNorm only, and both live. Configurations with $p$ use adaptive common-component scaling and the selected backward scale power; the remaining SplitCAM rows use the paper post-shifts. Fixed quantities are recomputed for each perturbed input. CAM baselines retain their library normalization and bilinear resize; SplitCAM retains signed maps and nearest-neighbor resize. All metrics apply their original absolute-value preprocessing. Gradient pairwise tests use archived rounded per-image scores; its original confidence intervals are retained. Three test illustrations were viewed during method development.
The paper-postshift search admitted no validation configuration for $L$, $U$.
